\begin{theorem}[有限 cyclotomic--Hurwitz 阴影对 \texorpdfstring{$S_4$}{S4} 椭圆谱宿主的不完备性]\label{thm:conclusion-finite-cyclotomic-hurwitz-shadows-incomplete-for-s4-elliptic-host}
设 \(\mathfrak H_{\pi}\) 为项目的谱宿主，即由正规化
\[
y:E\to \PP^1
\]
与四次谱方程 \(\Pi(\lambda,y)\) 所承载的一参数椭圆覆盖。若一个不变量 \(\mathcal I\) 仅通过有限多个下列数据因子化：
\begin{enumerate}
  \item 根值单位平均
  \[
  \frac1m\sum_{j=0}^{m-1}a_n(\omega_m^j),
  \]
  \item 周期 Dirichlet 阴影
  \[
  Z_A(s)=d_m^{-s}\sum_{r=1}^{d_m}\kappa_A(r)\,\zeta\!\left(s,\frac{r}{d_m}\right),
  \]
\end{enumerate}
则 \(\mathcal I\) 不可能成为 \(\mathfrak H_{\pi}\) 的完备不变量。更精确地，它不能确定项目谱曲线的正规化 \(y:E\to \PP^1\) 及其四次谱方程 \(\Pi(\lambda,y)\) 的 \(\QQ(y)\)-同构类，尤其不能恢复
\[
\operatorname{Gal}_{\QQ(y)}(\Pi(\lambda,y))\cong S_4
\]
所携带的四层单值化几何。
\end{theorem}
\begin{proof}
由定理 \ref{thm:conclusion-realinput40-root-unity-ghost-complete-primitive-degenerate} 与定理 \ref{thm:conclusion-periodic-dirichlet-shadow-hurwitz-finite-closure}，这里允许的全部阴影数据只依赖于根值单位平均、周期计数与 Hurwitz 残数账本。它们本质上都是 sheet-symmetric 的标量压缩：前者只保留 ghost 层的整除梳齿，后者只保留有限周期向量及其平均值。

相反，文稿已证明项目谱宿主的正规化是一条固定椭圆曲线 \(E\)，而
\[
T_m(y)=\operatorname{Tr}_{\QQ(E)/\QQ(y)}(\lambda^m)
\]
形成真正的四层单值化迹族；相应四次谱方程在 \(\QQ(y)\) 上具有满 \(S_4\) 伽罗瓦群，其分歧与片交换结构正是宿主的核心几何数据。任何完备不变量若要恢复 \(\Pi(\lambda,y)\) 的 \(\QQ(y)\)-同构类，就必须保留这种非阿贝尔片排列信息。

然而，上述有限 cyclotomic--Hurwitz 阴影先天抹去片交换信息，只保留对所有 sheet 对称的标量影子。因此它们不可能完备恢复 \(\mathfrak H_{\pi}\) 的 \(S_4\) 单值化几何。证毕。
\end{proof}

\begin{theorem}[Puiseux--Hurwitz 二分律]\label{thm:conclusion-puiseux-hurwitz-dichotomy}
设 \(I(\alpha)\) 为项目 Lee--Yang 权重体系的速率函数。则当 \(\alpha\downarrow 0\) 与 \(\alpha\uparrow \frac12\) 时，分别有
\[
I(\alpha)=\log 2+2\alpha\log\alpha+(\log 8-2)\alpha+O\!\bigl(\alpha^2\log(1/\alpha)\bigr),
\]
\[
I\!\left(\frac12-\delta\right)
=
\log 2+4\delta\log\delta+(12\log 2-4)\delta
+O\!\bigl(\delta^2\log(1/\delta)\bigr),
\qquad \delta\downarrow 0.
\]
其中系数对
\[
(2,4)
\]
分别来自 Perron 分支在 \(y\to 0^+\) 与 \(y\to +\infty\) 处的 Puiseux 指数 \(1/2\) 与 \(1/4\)。相反，定理 \ref{thm:conclusion-periodic-dirichlet-shadow-hurwitz-finite-closure} 的任意有限 Hurwitz 阴影只携带 \(s=1\) 的简单极点及其残数账本。故 \((2,4)\) 不是 Hurwitz 残数账本内部的量，而是椭圆谱宿主的分支几何不变量。
\end{theorem}
\begin{proof}
端点展开及其 Puiseux 来源均已由文稿中的 Lee--Yang 端点分析严格给出。另一方面，定理 \ref{thm:conclusion-periodic-dirichlet-shadow-hurwitz-finite-closure} 表明：任何有限周期 Dirichlet 阴影在标量层面都分解为有限个 Hurwitz \(\zeta\) 函数之和，其奇性机制只表现为 \(s=1\) 处的简单极点及有限残数账本。

二者的解析型别因此严格不同：端点展开中的 \(\alpha\log\alpha\) 与 \(\delta\log\delta\) 来自代数分支的 Puiseux 指数，而 Hurwitz 阴影只记录 meromorphic 极点的残数。故端点系数对 \((2,4)\) 不可能在 Hurwitz 残数账本内部生成，只能来自谱宿主的 Puiseux 几何。证毕。
\end{proof}

\begin{conjecture}[三寄存器分裂原理]\label{conj:conclusion-ghost-hurwitz-sheet-three-register-splitting}
任一可精确嵌入项目的 \(\pi\)-通道，都应存在函子性分裂
\[
\mathfrak R_{\mathrm{ghost}}
\oplus
\mathfrak R_{\mathrm{Hurwitz}}
\oplus
\mathfrak R_{\mathrm{sheet}},
\]
其中
\[
\mathfrak R_{\mathrm{ghost}}
\quad\text{负责定理 \ref{thm:conclusion-realinput40-root-unity-ghost-complete-primitive-degenerate} 的 Dirichlet 梳齿层析，}
\]
\[
\mathfrak R_{\mathrm{Hurwitz}}
\quad\text{负责定理 \ref{thm:conclusion-periodic-dirichlet-shadow-hurwitz-finite-closure} 的有限周期残数账本，}
\]
\[
\mathfrak R_{\mathrm{sheet}}
\quad\text{负责定理 \ref{thm:conclusion-finite-cyclotomic-hurwitz-shadows-incomplete-for-s4-elliptic-host} 与定理 \ref{thm:conclusion-puiseux-hurwitz-dichotomy} 中的椭圆 \(S_4\) 单值化及端点 Puiseux 奇性。}
\]
并且，缺失第三寄存器的任何公式族都只能给出宿主的投影读出，而不可能与宿主本身同构；Machin-like 与 BBP-like 阴影应当都属于这一情形。
\end{conjecture}
